
\documentclass{article}
\usepackage{colortbl}
\usepackage{makecell}
\usepackage{multirow}
\usepackage{supertabular}

\begin{document}

\newcounter{utterance}

\centering \large Interaction Transcript for game `dond', experiment `coop\_en', episode 8 with Qwen3{-}32B{-}t0.0.
\vspace{24pt}

{ \footnotesize  \setcounter{utterance}{1}
\setlength{\tabcolsep}{0pt}
\begin{supertabular}{c@{$\;$}|p{.15\linewidth}@{}p{.15\linewidth}p{.15\linewidth}p{.15\linewidth}p{.15\linewidth}p{.15\linewidth}}
    \# & $\;$A & \multicolumn{4}{c}{Game Master} & $\;\:$B\\
    \hline

    \theutterance \stepcounter{utterance}  
    & & \multicolumn{4}{p{0.6\linewidth}}{
        \cellcolor[rgb]{0.9,0.9,0.9}{
            \makecell[{{p{\linewidth}}}]{
                \texttt{\tiny{[P1$\langle$GM]}}
                \texttt{You are playing a cooperative negotiation game in which you have to agree on how to divide a set of items among you and another player.} \\
\\ 
\texttt{Rules:} \\
\texttt{(a) You and the other player are given a set of items. Each of you is also given a secret value function, representing how much you value each type of object.} \\
\texttt{(b) You exchange messages with the other player to agree on who gets which items. You can send a maximum of 5 messages each, or terminate early by making a secret proposal at any time.} \\
\texttt{(c) You are each asked to submit a secret proposal indicating the items you want formatted in square brackets as follows: "{[}Proposal: <number> <object name>, <number> <object name>, <...>{]}"} \\
\texttt{(d) If your proposals are complementary, i.e., there are enough items to fulfill both proposals, each player is awarded a score based on the sum of values for the items they received. Otherwise, both of you get zero points.} \\
\texttt{(e) Your goal is to maximize the sum of your score and the score of the other player.} \\
\\ 
\texttt{Let us start.} \\
\\ 
\texttt{The set of available items is:} \\
\\ 
\texttt{3 scissors, 1 thimble, 1 fork} \\
\\ 
\texttt{Your secret value function is:} \\
\\ 
\texttt{scissors: 3, thimble: 0, fork: 1} \\
\\ 
\texttt{IMPORTANT: Your messages, unless it is the secret proposal, are directly transmitted to the other player, so do not include any response to the rules or text announcing your message. To make a secret proposal, use the indicated format. Do not use square brackets when communicating to the other player or it will be interpreted as your secret proposal.} \\
            }
        }
    }
    & \\ \\

    \theutterance \stepcounter{utterance}  
    & \multicolumn{4}{p{0.6\linewidth}}{
        \cellcolor[rgb]{0.9,0.9,0.9}{
            \makecell[{{p{\linewidth}}}]{
                \texttt{\tiny{[P1$\rangle$GM]}}
                \texttt{<think>} \\
\texttt{Okay, let's see. I need to divide 3 scissors, 1 thimble, and 1 fork between me and another player. My value function is scissors: 3, thimble: 0, fork: 1. The other player's values are secret, so I don't know what they value. The goal is to maximize the total score for both of us. } \\
\\ 
\texttt{First, I should figure out what the other player might value. Since I don't have their value function, I need to communicate to find out. But I can only send 5 messages. Let me think about possible strategies.} \\
\\ 
\texttt{The items are 3 scissors, 1 thimble, 1 fork. My highest value is scissors, then fork, then thimble. The other player might value the thimble or fork more. Since there's only one thimble and one fork, maybe they want those. But there are three scissors, so maybe they can take some scissors as well.} \\
\\ 
\texttt{If I propose taking all 3 scissors and the fork, that would give me 3*3 +1 = 10 points. But then the other player would have to take the thimble and maybe some scissors. But if their value for thimble is high, that might work. However, if they also value scissors, maybe they want some. But since there are three, maybe we can split them. } \\
\\ 
\texttt{Alternatively, if I take 2 scissors and the fork, that's 2*3 +1 =7, leaving 1 scissors and thimble for the other. But if they value the thimble more, that could be better. But I don't know their values. } \\
\\ 
\texttt{Wait, the goal is to maximize the total score. So maybe we should allocate items where each player gets the items they value most. Since I don't know their values, I need to ask. But how? Since I can't ask directly, maybe I can make a proposal and see if they accept or not. But since we can make secret proposals, maybe I can propose a split and see if it's complementary. } \\
\\ 
\texttt{But the problem is that if I make a secret proposal, the other player has to make one too. If our proposals are complementary, then we get our scores. Otherwise, zero. So I need to make a proposal that is likely to be complementary with theirs. } \\
\\ 
\texttt{Since I don't know their values} \\
            }
        }
    }
    & & \\ \\

    \theutterance \stepcounter{utterance}  
    & & & \multicolumn{2}{p{0.3\linewidth}}{
        \cellcolor[rgb]{0.9,0.9,0.9}{
            \makecell[{{p{\linewidth}}}]{
                \texttt{\tiny{[GM$|$GM]}}
                \texttt{continue} \\
            }
        }
    }
    & & \\ \\

    \theutterance \stepcounter{utterance}  
    & & & \multicolumn{2}{p{0.3\linewidth}}{
        \cellcolor[rgb]{0.9,0.9,0.9}{
            \makecell[{{p{\linewidth}}}]{
                \texttt{\tiny{[GM$|$GM]}}
                \texttt{<think>} \\
\texttt{Okay, let's see. I need to divide 3 scissors, 1 thimble, and 1 fork between me and another player. My value function is scissors: 3, thimble: 0, fork: 1. The other player's values are secret, so I don't know what they value. The goal is to maximize the total score for both of us. } \\
\\ 
\texttt{First, I should figure out what the other player might value. Since I don't have their value function, I need to communicate to find out. But I can only send 5 messages. Let me think about possible strategies.} \\
\\ 
\texttt{The items are 3 scissors, 1 thimble, 1 fork. My highest value is scissors, then fork, then thimble. The other player might value the thimble or fork more. Since there's only one thimble and one fork, maybe they want those. But there are three scissors, so maybe they can take some scissors as well.} \\
\\ 
\texttt{If I propose taking all 3 scissors and the fork, that would give me 3*3 +1 = 10 points. But then the other player would have to take the thimble and maybe some scissors. But if their value for thimble is high, that might work. However, if they also value scissors, maybe they want some. But since there are three, maybe we can split them. } \\
\\ 
\texttt{Alternatively, if I take 2 scissors and the fork, that's 2*3 +1 =7, leaving 1 scissors and thimble for the other. But if they value the thimble more, that could be better. But I don't know their values. } \\
\\ 
\texttt{Wait, the goal is to maximize the total score. So maybe we should allocate items where each player gets the items they value most. Since I don't know their values, I need to ask. But how? Since I can't ask directly, maybe I can make a proposal and see if they accept or not. But since we can make secret proposals, maybe I can propose a split and see if it's complementary. } \\
\\ 
\texttt{But the problem is that if I make a secret proposal, the other player has to make one too. If our proposals are complementary, then we get our scores. Otherwise, zero. So I need to make a proposal that is likely to be complementary with theirs. } \\
\\ 
\texttt{Since I don't know their values} \\
            }
        }
    }
    & & \\ \\

    \theutterance \stepcounter{utterance}  
    & & \multicolumn{4}{p{0.6\linewidth}}{
        \cellcolor[rgb]{0.9,0.9,0.9}{
            \makecell[{{p{\linewidth}}}]{
                \texttt{\tiny{[GM$\rangle$P2]}}
                \texttt{You are playing a cooperative negotiation game in which you have to agree on how to divide a set of items among you and another player.} \\
\\ 
\texttt{Rules:} \\
\texttt{(a) You and the other player are given a set of items. Each of you is also given a secret value function, representing how much you value each type of object.} \\
\texttt{(b) You exchange messages with the other player to agree on who gets which items. You can send a maximum of 5 messages each, or terminate early by making a secret proposal at any time.} \\
\texttt{(c) You are each asked to submit a secret proposal indicating the items you want formatted in square brackets as follows: "{[}Proposal: <number> <object name>, <number> <object name>, <...>{]}"} \\
\texttt{(d) If your proposals are complementary, i.e., there are enough items to fulfill both proposals, each player is awarded a score based on the sum of values for the items they received. Otherwise, both of you get zero points.} \\
\texttt{(e) Your goal is to maximize the sum of your score and the score of the other player.} \\
\\ 
\texttt{Let us start.} \\
\\ 
\texttt{The set of available items is:} \\
\\ 
\texttt{3 scissors, 1 thimble, 1 fork} \\
\\ 
\texttt{Your secret value function is:} \\
\\ 
\texttt{scissors: 1, thimble: 6, fork: 1} \\
\\ 
\texttt{IMPORTANT: Your messages, unless it is the secret proposal, are directly transmitted to the other player, so do not include any response to the rules or text announcing your message. To make a secret proposal, use the indicated format. Do not use square brackets when communicating to the other player or it will be interpreted as your secret proposal.} \\
\\ 
\\ 
\texttt{<think>} \\
\texttt{Okay, let's see. I need to divide 3 scissors, 1 thimble, and 1 fork between me and another player. My value function is scissors: 3, thimble: 0, fork: 1. The other player's values are secret, so I don't know what they value. The goal is to maximize the total score for both of us. } \\
\\ 
\texttt{First, I should figure out what the other player might value. Since I don't have their value function, I need to communicate to find out. But I can only send 5 messages. Let me think about possible strategies.} \\
\\ 
\texttt{The items are 3 scissors, 1 thimble, 1 fork. My highest value is scissors, then fork, then thimble. The other player might value the thimble or fork more. Since there's only one thimble and one fork, maybe they want those. But there are three scissors, so maybe they can take some scissors as well.} \\
\\ 
\texttt{If I propose taking all 3 scissors and the fork, that would give me 3*3 +1 = 10 points. But then the other player would have to take the thimble and maybe some scissors. But if their value for thimble is high, that might work. However, if they also value scissors, maybe they want some. But since there are three, maybe we can split them. } \\
\\ 
\texttt{Alternatively, if I take 2 scissors and the fork, that's 2*3 +1 =7, leaving 1 scissors and thimble for the other. But if they value the thimble more, that could be better. But I don't know their values. } \\
\\ 
\texttt{Wait, the goal is to maximize the total score. So maybe we should allocate items where each player gets the items they value most. Since I don't know their values, I need to ask. But how? Since I can't ask directly, maybe I can make a proposal and see if they accept or not. But since we can make secret proposals, maybe I can propose a split and see if it's complementary. } \\
\\ 
\texttt{But the problem is that if I make a secret proposal, the other player has to make one too. If our proposals are complementary, then we get our scores. Otherwise, zero. So I need to make a proposal that is likely to be complementary with theirs. } \\
\\ 
\texttt{Since I don't know their values} \\
            }
        }
    }
    & \\ \\

    \theutterance \stepcounter{utterance}  
    & & & \multicolumn{4}{p{0.6\linewidth}}{
        \cellcolor[rgb]{0.9,0.9,0.9}{
            \makecell[{{p{\linewidth}}}]{
                \texttt{\tiny{[GM$\langle$P2]}}
                \texttt{<think>} \\
\texttt{Okay, let's try to figure this out. The user is playing a negotiation game where they need to divide items between themselves and another player. The items are 3 scissors, 1 thimble, and 1 fork. The user's value function is scissors: 3, thimble: 0, fork: 1. The other player's values are secret. The goal is to maximize the total score for both players.} \\
\\ 
\texttt{First, the user needs to consider possible value functions for the other player. Since the thimble and fork are unique, the other player might value them highly. The scissors are three in number, so maybe the other player values them less, but it's possible they value them more. The user's highest value is on scissors, so they might want as many as possible, but they also have a fork worth 1. The thimble is worth 0 to the user, so ideally, the other player should get it if they value it.} \\
\\ 
\texttt{The challenge is to communicate effectively within 5 messages to find a complementary proposal. The user can't know the other's values, so they need to make a proposal that allows for a good split. Since the user's highest value is scissors, maybe they should take as many as possible, but leave some for the other player. However, the fork and thimble need to be allocated. If the user takes the fork, the other player might need the thimble. Alternatively, if the other player values the fork more, the user could take the thimble (even though it's worth 0 to them) to let the other have the fork.} \\
\\ 
\texttt{But how to figure out the other's preferences? The user could propose different splits and see if the other's proposal is complementary. For example, the user could propose taking 2 scissors and the fork, leaving 1 scissors and the thimble. If the other player values the thimble, that might work. Alternatively, if the other player wants the fork, the user could take the thimble and some scissors.} \\
\\ 
\texttt{Wait, the user's value for the thimble is 0, so it's better to let the other player have it if possible. The fork is worth 1 to the user, so maybe the user should take it if the other player doesn't value it highly. But without knowing the other's values, it's tricky.} \\
\\ 
\texttt{A possible strategy is to propose taking} \\
            }
        }
    }
     \\ \\

    \theutterance \stepcounter{utterance}  
    & & & \multicolumn{2}{p{0.3\linewidth}}{
        \cellcolor[rgb]{0.9,0.9,0.9}{
            \makecell[{{p{\linewidth}}}]{
                \texttt{\tiny{[GM$|$GM]}}
                \texttt{continue} \\
            }
        }
    }
    & & \\ \\

    \theutterance \stepcounter{utterance}  
    & & & \multicolumn{2}{p{0.3\linewidth}}{
        \cellcolor[rgb]{0.9,0.9,0.9}{
            \makecell[{{p{\linewidth}}}]{
                \texttt{\tiny{[GM$|$GM]}}
                \texttt{<think>} \\
\texttt{Okay, let's try to figure this out. The user is playing a negotiation game where they need to divide items between themselves and another player. The items are 3 scissors, 1 thimble, and 1 fork. The user's value function is scissors: 3, thimble: 0, fork: 1. The other player's values are secret. The goal is to maximize the total score for both players.} \\
\\ 
\texttt{First, the user needs to consider possible value functions for the other player. Since the thimble and fork are unique, the other player might value them highly. The scissors are three in number, so maybe the other player values them less, but it's possible they value them more. The user's highest value is on scissors, so they might want as many as possible, but they also have a fork worth 1. The thimble is worth 0 to the user, so ideally, the other player should get it if they value it.} \\
\\ 
\texttt{The challenge is to communicate effectively within 5 messages to find a complementary proposal. The user can't know the other's values, so they need to make a proposal that allows for a good split. Since the user's highest value is scissors, maybe they should take as many as possible, but leave some for the other player. However, the fork and thimble need to be allocated. If the user takes the fork, the other player might need the thimble. Alternatively, if the other player values the fork more, the user could take the thimble (even though it's worth 0 to them) to let the other have the fork.} \\
\\ 
\texttt{But how to figure out the other's preferences? The user could propose different splits and see if the other's proposal is complementary. For example, the user could propose taking 2 scissors and the fork, leaving 1 scissors and the thimble. If the other player values the thimble, that might work. Alternatively, if the other player wants the fork, the user could take the thimble and some scissors.} \\
\\ 
\texttt{Wait, the user's value for the thimble is 0, so it's better to let the other player have it if possible. The fork is worth 1 to the user, so maybe the user should take it if the other player doesn't value it highly. But without knowing the other's values, it's tricky.} \\
\\ 
\texttt{A possible strategy is to propose taking} \\
            }
        }
    }
    & & \\ \\

    \theutterance \stepcounter{utterance}  
    & & \multicolumn{4}{p{0.6\linewidth}}{
        \cellcolor[rgb]{0.9,0.9,0.9}{
            \makecell[{{p{\linewidth}}}]{
                \texttt{\tiny{[P1$\langle$GM]}}
                \texttt{<think>} \\
\texttt{Okay, let's try to figure this out. The user is playing a negotiation game where they need to divide items between themselves and another player. The items are 3 scissors, 1 thimble, and 1 fork. The user's value function is scissors: 3, thimble: 0, fork: 1. The other player's values are secret. The goal is to maximize the total score for both players.} \\
\\ 
\texttt{First, the user needs to consider possible value functions for the other player. Since the thimble and fork are unique, the other player might value them highly. The scissors are three in number, so maybe the other player values them less, but it's possible they value them more. The user's highest value is on scissors, so they might want as many as possible, but they also have a fork worth 1. The thimble is worth 0 to the user, so ideally, the other player should get it if they value it.} \\
\\ 
\texttt{The challenge is to communicate effectively within 5 messages to find a complementary proposal. The user can't know the other's values, so they need to make a proposal that allows for a good split. Since the user's highest value is scissors, maybe they should take as many as possible, but leave some for the other player. However, the fork and thimble need to be allocated. If the user takes the fork, the other player might need the thimble. Alternatively, if the other player values the fork more, the user could take the thimble (even though it's worth 0 to them) to let the other have the fork.} \\
\\ 
\texttt{But how to figure out the other's preferences? The user could propose different splits and see if the other's proposal is complementary. For example, the user could propose taking 2 scissors and the fork, leaving 1 scissors and the thimble. If the other player values the thimble, that might work. Alternatively, if the other player wants the fork, the user could take the thimble and some scissors.} \\
\\ 
\texttt{Wait, the user's value for the thimble is 0, so it's better to let the other player have it if possible. The fork is worth 1 to the user, so maybe the user should take it if the other player doesn't value it highly. But without knowing the other's values, it's tricky.} \\
\\ 
\texttt{A possible strategy is to propose taking} \\
            }
        }
    }
    & \\ \\

    \theutterance \stepcounter{utterance}  
    & \multicolumn{4}{p{0.6\linewidth}}{
        \cellcolor[rgb]{0.9,0.9,0.9}{
            \makecell[{{p{\linewidth}}}]{
                \texttt{\tiny{[P1$\rangle$GM]}}
                \texttt{<think>} \\
\texttt{Okay, let's try to figure out the best way to approach this. The user has to divide 3 scissors, 1 thimble, and 1 fork between themselves and another player. Their own values are scissors: 3, thimble: 0, fork: 1. The other player's values are secret, so the user needs to negotiate to maximize the total score.} \\
\\ 
\texttt{First, the user's highest value is on scissors, so they probably want as many as possible. But there are three scissors, so maybe the other player can take some too. The fork is worth 1 to the user, and the thimble is worth 0. The other player might value the thimble or the fork more. Since there's only one of each, it's important to allocate them where they're most valued.} \\
\\ 
\texttt{The user can send up to 5 messages. They need to communicate their preferences and find out the other player's. But since the other player's values are secret, the user has to make proposals and see if they're complementary.} \\
\\ 
\texttt{One approach is to propose taking 2 scissors and the fork, leaving 1 scissors and the thimble. That way, the user gets 2*3 +1 =7 points. If the other player values the thimble, they can take it and the remaining scissors. But if the other player values the fork more, maybe they'd want it instead. Alternatively, the user could take 1 scissors and the fork, leaving 2 scissors and the thimble. That would give the user 3 +1 =4, but maybe the other player can get more from the remaining items.} \\
\\ 
\texttt{Wait, the goal is to maximize the total score. So the user should aim for a split where both players get the items they value most. Since the user's thimble is 0, it's better to let the other player have it if they value it. The fork is worth 1 to the user, so if the other player values it more, it's better to let them have it. But how to find out?} \\
\\ 
\texttt{Maybe the user should propose different splits and see if the other player's proposal is complementary. For example, if the user proposes taking 2 scissors and the fork, and the other player takes 1 scissors and the thimble, that's complementary. The user's score would be 7, and the other's depends on their values} \\
            }
        }
    }
    & & \\ \\

    \theutterance \stepcounter{utterance}  
    & & & \multicolumn{2}{p{0.3\linewidth}}{
        \cellcolor[rgb]{0.9,0.9,0.9}{
            \makecell[{{p{\linewidth}}}]{
                \texttt{\tiny{[GM$|$GM]}}
                \texttt{continue} \\
            }
        }
    }
    & & \\ \\

    \theutterance \stepcounter{utterance}  
    & & & \multicolumn{2}{p{0.3\linewidth}}{
        \cellcolor[rgb]{0.9,0.9,0.9}{
            \makecell[{{p{\linewidth}}}]{
                \texttt{\tiny{[GM$|$GM]}}
                \texttt{<think>} \\
\texttt{Okay, let's try to figure out the best way to approach this. The user has to divide 3 scissors, 1 thimble, and 1 fork between themselves and another player. Their own values are scissors: 3, thimble: 0, fork: 1. The other player's values are secret, so the user needs to negotiate to maximize the total score.} \\
\\ 
\texttt{First, the user's highest value is on scissors, so they probably want as many as possible. But there are three scissors, so maybe the other player can take some too. The fork is worth 1 to the user, and the thimble is worth 0. The other player might value the thimble or the fork more. Since there's only one of each, it's important to allocate them where they're most valued.} \\
\\ 
\texttt{The user can send up to 5 messages. They need to communicate their preferences and find out the other player's. But since the other player's values are secret, the user has to make proposals and see if they're complementary.} \\
\\ 
\texttt{One approach is to propose taking 2 scissors and the fork, leaving 1 scissors and the thimble. That way, the user gets 2*3 +1 =7 points. If the other player values the thimble, they can take it and the remaining scissors. But if the other player values the fork more, maybe they'd want it instead. Alternatively, the user could take 1 scissors and the fork, leaving 2 scissors and the thimble. That would give the user 3 +1 =4, but maybe the other player can get more from the remaining items.} \\
\\ 
\texttt{Wait, the goal is to maximize the total score. So the user should aim for a split where both players get the items they value most. Since the user's thimble is 0, it's better to let the other player have it if they value it. The fork is worth 1 to the user, so if the other player values it more, it's better to let them have it. But how to find out?} \\
\\ 
\texttt{Maybe the user should propose different splits and see if the other player's proposal is complementary. For example, if the user proposes taking 2 scissors and the fork, and the other player takes 1 scissors and the thimble, that's complementary. The user's score would be 7, and the other's depends on their values} \\
            }
        }
    }
    & & \\ \\

    \theutterance \stepcounter{utterance}  
    & & \multicolumn{4}{p{0.6\linewidth}}{
        \cellcolor[rgb]{0.9,0.9,0.9}{
            \makecell[{{p{\linewidth}}}]{
                \texttt{\tiny{[GM$\rangle$P2]}}
                \texttt{<think>} \\
\texttt{Okay, let's try to figure out the best way to approach this. The user has to divide 3 scissors, 1 thimble, and 1 fork between themselves and another player. Their own values are scissors: 3, thimble: 0, fork: 1. The other player's values are secret, so the user needs to negotiate to maximize the total score.} \\
\\ 
\texttt{First, the user's highest value is on scissors, so they probably want as many as possible. But there are three scissors, so maybe the other player can take some too. The fork is worth 1 to the user, and the thimble is worth 0. The other player might value the thimble or the fork more. Since there's only one of each, it's important to allocate them where they're most valued.} \\
\\ 
\texttt{The user can send up to 5 messages. They need to communicate their preferences and find out the other player's. But since the other player's values are secret, the user has to make proposals and see if they're complementary.} \\
\\ 
\texttt{One approach is to propose taking 2 scissors and the fork, leaving 1 scissors and the thimble. That way, the user gets 2*3 +1 =7 points. If the other player values the thimble, they can take it and the remaining scissors. But if the other player values the fork more, maybe they'd want it instead. Alternatively, the user could take 1 scissors and the fork, leaving 2 scissors and the thimble. That would give the user 3 +1 =4, but maybe the other player can get more from the remaining items.} \\
\\ 
\texttt{Wait, the goal is to maximize the total score. So the user should aim for a split where both players get the items they value most. Since the user's thimble is 0, it's better to let the other player have it if they value it. The fork is worth 1 to the user, so if the other player values it more, it's better to let them have it. But how to find out?} \\
\\ 
\texttt{Maybe the user should propose different splits and see if the other player's proposal is complementary. For example, if the user proposes taking 2 scissors and the fork, and the other player takes 1 scissors and the thimble, that's complementary. The user's score would be 7, and the other's depends on their values} \\
            }
        }
    }
    & \\ \\

    \theutterance \stepcounter{utterance}  
    & & & \multicolumn{4}{p{0.6\linewidth}}{
        \cellcolor[rgb]{0.9,0.9,0.9}{
            \makecell[{{p{\linewidth}}}]{
                \texttt{\tiny{[GM$\langle$P2]}}
                \texttt{<think>} \\
\texttt{Okay, let's try to figure out the best way to approach this. The user has to divide 3 scissors, 1 thimble, and 1 fork between themselves and another player. Their own values are scissors: 3, thimble: 0, fork: 1. The other player's values are secret, so the user needs to negotiate to maximize the total score.} \\
\\ 
\texttt{First, the user's highest value is on scissors, so they probably want as many as possible. But there are three scissors, so maybe the other player can take some too. The fork is worth 1 to the user, and the thimble is worth 0. The other player might value the thimble or the fork more. Since there's only one of each, it's important to allocate them where they're most valued.} \\
\\ 
\texttt{The user can send up to 5 messages. They need to communicate their preferences and find out the other player's. But since the other player's values are secret, the user has to make proposals and see if they're complementary.} \\
\\ 
\texttt{One approach is to propose taking 2 scissors and the fork, leaving 1 scissors and the thimble. That way, the user gets 2*3 +1 =7 points. If the other player values the thimble, they can take it and the remaining scissors. But if the other player values the fork more, maybe they'd want it instead. Alternatively, the user could take 1 scissors and the fork, leaving 2 scissors and the thimble. That would give the user 3 +1 =4, but maybe the other player can get more from the remaining items.} \\
\\ 
\texttt{Wait, the goal is to maximize the total score. So the user should aim for a split where both players get the items they value most. Since the user's thimble is 0, it's better to let the other player have it if they value it. The fork is worth 1 to the user, so if the other player values it more, it's better to let them have it. But how to find out?} \\
\\ 
\texttt{Maybe the user should propose different splits and see if the other player's proposal is complementary. For example, if the user proposes taking 2 scissors and the fork, and the other player takes 1 scissors and the thimble, that's complementary. The user's score would be 7, and the other's depends on their values} \\
            }
        }
    }
     \\ \\

    \theutterance \stepcounter{utterance}  
    & & & \multicolumn{2}{p{0.3\linewidth}}{
        \cellcolor[rgb]{0.9,0.9,0.9}{
            \makecell[{{p{\linewidth}}}]{
                \texttt{\tiny{[GM$|$GM]}}
                \texttt{continue} \\
            }
        }
    }
    & & \\ \\

    \theutterance \stepcounter{utterance}  
    & & & \multicolumn{2}{p{0.3\linewidth}}{
        \cellcolor[rgb]{0.9,0.9,0.9}{
            \makecell[{{p{\linewidth}}}]{
                \texttt{\tiny{[GM$|$GM]}}
                \texttt{<think>} \\
\texttt{Okay, let's try to figure out the best way to approach this. The user has to divide 3 scissors, 1 thimble, and 1 fork between themselves and another player. Their own values are scissors: 3, thimble: 0, fork: 1. The other player's values are secret, so the user needs to negotiate to maximize the total score.} \\
\\ 
\texttt{First, the user's highest value is on scissors, so they probably want as many as possible. But there are three scissors, so maybe the other player can take some too. The fork is worth 1 to the user, and the thimble is worth 0. The other player might value the thimble or the fork more. Since there's only one of each, it's important to allocate them where they're most valued.} \\
\\ 
\texttt{The user can send up to 5 messages. They need to communicate their preferences and find out the other player's. But since the other player's values are secret, the user has to make proposals and see if they're complementary.} \\
\\ 
\texttt{One approach is to propose taking 2 scissors and the fork, leaving 1 scissors and the thimble. That way, the user gets 2*3 +1 =7 points. If the other player values the thimble, they can take it and the remaining scissors. But if the other player values the fork more, maybe they'd want it instead. Alternatively, the user could take 1 scissors and the fork, leaving 2 scissors and the thimble. That would give the user 3 +1 =4, but maybe the other player can get more from the remaining items.} \\
\\ 
\texttt{Wait, the goal is to maximize the total score. So the user should aim for a split where both players get the items they value most. Since the user's thimble is 0, it's better to let the other player have it if they value it. The fork is worth 1 to the user, so if the other player values it more, it's better to let them have it. But how to find out?} \\
\\ 
\texttt{Maybe the user should propose different splits and see if the other player's proposal is complementary. For example, if the user proposes taking 2 scissors and the fork, and the other player takes 1 scissors and the thimble, that's complementary. The user's score would be 7, and the other's depends on their values} \\
            }
        }
    }
    & & \\ \\

    \theutterance \stepcounter{utterance}  
    & & \multicolumn{4}{p{0.6\linewidth}}{
        \cellcolor[rgb]{0.9,0.9,0.9}{
            \makecell[{{p{\linewidth}}}]{
                \texttt{\tiny{[P1$\langle$GM]}}
                \texttt{<think>} \\
\texttt{Okay, let's try to figure out the best way to approach this. The user has to divide 3 scissors, 1 thimble, and 1 fork between themselves and another player. Their own values are scissors: 3, thimble: 0, fork: 1. The other player's values are secret, so the user needs to negotiate to maximize the total score.} \\
\\ 
\texttt{First, the user's highest value is on scissors, so they probably want as many as possible. But there are three scissors, so maybe the other player can take some too. The fork is worth 1 to the user, and the thimble is worth 0. The other player might value the thimble or the fork more. Since there's only one of each, it's important to allocate them where they're most valued.} \\
\\ 
\texttt{The user can send up to 5 messages. They need to communicate their preferences and find out the other player's. But since the other player's values are secret, the user has to make proposals and see if they're complementary.} \\
\\ 
\texttt{One approach is to propose taking 2 scissors and the fork, leaving 1 scissors and the thimble. That way, the user gets 2*3 +1 =7 points. If the other player values the thimble, they can take it and the remaining scissors. But if the other player values the fork more, maybe they'd want it instead. Alternatively, the user could take 1 scissors and the fork, leaving 2 scissors and the thimble. That would give the user 3 +1 =4, but maybe the other player can get more from the remaining items.} \\
\\ 
\texttt{Wait, the goal is to maximize the total score. So the user should aim for a split where both players get the items they value most. Since the user's thimble is 0, it's better to let the other player have it if they value it. The fork is worth 1 to the user, so if the other player values it more, it's better to let them have it. But how to find out?} \\
\\ 
\texttt{Maybe the user should propose different splits and see if the other player's proposal is complementary. For example, if the user proposes taking 2 scissors and the fork, and the other player takes 1 scissors and the thimble, that's complementary. The user's score would be 7, and the other's depends on their values} \\
            }
        }
    }
    & \\ \\

    \theutterance \stepcounter{utterance}  
    & \multicolumn{4}{p{0.6\linewidth}}{
        \cellcolor[rgb]{0.9,0.9,0.9}{
            \makecell[{{p{\linewidth}}}]{
                \texttt{\tiny{[P1$\rangle$GM]}}
                \texttt{<think>} \\
\texttt{Okay, let's try to figure out the best way to approach this. The user has to divide 3 scissors, 1 thimble, and 1 fork between themselves and another player. Their own values are scissors: 3, thimble: 0, fork: 1. The other player's values are secret, so the user needs to negotiate to maximize the total score.} \\
\\ 
\texttt{First, the user's highest value is on scissors, so they probably want as many as possible. But there are three scissors, so maybe the other player can take some too. The fork is worth 1 to the user, and the thimble is worth 0. The other player might value the thimble or the fork more. Since there's only one of each, it's important to allocate them where they're most valued.} \\
\\ 
\texttt{The user can send up to 5 messages. They need to communicate their preferences and find out the other player's. But since the other player's values are secret, the user has to make proposals and see if they're complementary.} \\
\\ 
\texttt{One approach is to propose taking 2 scissors and the fork, leaving 1 scissors and the thimble. That way, the user gets 2*3 +1 =7 points. If the other player values the thimble, they can take it and the remaining scissors. But if the other player values the fork more, maybe they'd want it instead. Alternatively, the user could take 1 scissors and the fork, leaving 2 scissors and the thimble. That would give the user 3 +1 =4, but maybe the other player can get more from the remaining items.} \\
\\ 
\texttt{Wait, the goal is to maximize the total score. So the user should aim for a split where both players get the items they value most. Since the user's thimble is 0, it's better to let the other player have it if they value it. The fork is worth 1 to the user, so if the other player values it more, it's better to let them have it. But how to find out?} \\
\\ 
\texttt{Maybe the user should propose different splits and see if the other player's proposal is complementary. For example, if the user proposes taking 2 scissors and the fork, and the other player takes 1 scissors and the thimble, that's complementary. The user's score would be 7, and the other's depends on their values} \\
            }
        }
    }
    & & \\ \\

    \theutterance \stepcounter{utterance}  
    & & & \multicolumn{2}{p{0.3\linewidth}}{
        \cellcolor[rgb]{0.9,0.9,0.9}{
            \makecell[{{p{\linewidth}}}]{
                \texttt{\tiny{[GM$|$GM]}}
                \texttt{continue} \\
            }
        }
    }
    & & \\ \\

    \theutterance \stepcounter{utterance}  
    & & & \multicolumn{2}{p{0.3\linewidth}}{
        \cellcolor[rgb]{0.9,0.9,0.9}{
            \makecell[{{p{\linewidth}}}]{
                \texttt{\tiny{[GM$|$GM]}}
                \texttt{<think>} \\
\texttt{Okay, let's try to figure out the best way to approach this. The user has to divide 3 scissors, 1 thimble, and 1 fork between themselves and another player. Their own values are scissors: 3, thimble: 0, fork: 1. The other player's values are secret, so the user needs to negotiate to maximize the total score.} \\
\\ 
\texttt{First, the user's highest value is on scissors, so they probably want as many as possible. But there are three scissors, so maybe the other player can take some too. The fork is worth 1 to the user, and the thimble is worth 0. The other player might value the thimble or the fork more. Since there's only one of each, it's important to allocate them where they're most valued.} \\
\\ 
\texttt{The user can send up to 5 messages. They need to communicate their preferences and find out the other player's. But since the other player's values are secret, the user has to make proposals and see if they're complementary.} \\
\\ 
\texttt{One approach is to propose taking 2 scissors and the fork, leaving 1 scissors and the thimble. That way, the user gets 2*3 +1 =7 points. If the other player values the thimble, they can take it and the remaining scissors. But if the other player values the fork more, maybe they'd want it instead. Alternatively, the user could take 1 scissors and the fork, leaving 2 scissors and the thimble. That would give the user 3 +1 =4, but maybe the other player can get more from the remaining items.} \\
\\ 
\texttt{Wait, the goal is to maximize the total score. So the user should aim for a split where both players get the items they value most. Since the user's thimble is 0, it's better to let the other player have it if they value it. The fork is worth 1 to the user, so if the other player values it more, it's better to let them have it. But how to find out?} \\
\\ 
\texttt{Maybe the user should propose different splits and see if the other player's proposal is complementary. For example, if the user proposes taking 2 scissors and the fork, and the other player takes 1 scissors and the thimble, that's complementary. The user's score would be 7, and the other's depends on their values} \\
            }
        }
    }
    & & \\ \\

    \theutterance \stepcounter{utterance}  
    & & \multicolumn{4}{p{0.6\linewidth}}{
        \cellcolor[rgb]{0.9,0.9,0.9}{
            \makecell[{{p{\linewidth}}}]{
                \texttt{\tiny{[GM$\rangle$P2]}}
                \texttt{<think>} \\
\texttt{Okay, let's try to figure out the best way to approach this. The user has to divide 3 scissors, 1 thimble, and 1 fork between themselves and another player. Their own values are scissors: 3, thimble: 0, fork: 1. The other player's values are secret, so the user needs to negotiate to maximize the total score.} \\
\\ 
\texttt{First, the user's highest value is on scissors, so they probably want as many as possible. But there are three scissors, so maybe the other player can take some too. The fork is worth 1 to the user, and the thimble is worth 0. The other player might value the thimble or the fork more. Since there's only one of each, it's important to allocate them where they're most valued.} \\
\\ 
\texttt{The user can send up to 5 messages. They need to communicate their preferences and find out the other player's. But since the other player's values are secret, the user has to make proposals and see if they're complementary.} \\
\\ 
\texttt{One approach is to propose taking 2 scissors and the fork, leaving 1 scissors and the thimble. That way, the user gets 2*3 +1 =7 points. If the other player values the thimble, they can take it and the remaining scissors. But if the other player values the fork more, maybe they'd want it instead. Alternatively, the user could take 1 scissors and the fork, leaving 2 scissors and the thimble. That would give the user 3 +1 =4, but maybe the other player can get more from the remaining items.} \\
\\ 
\texttt{Wait, the goal is to maximize the total score. So the user should aim for a split where both players get the items they value most. Since the user's thimble is 0, it's better to let the other player have it if they value it. The fork is worth 1 to the user, so if the other player values it more, it's better to let them have it. But how to find out?} \\
\\ 
\texttt{Maybe the user should propose different splits and see if the other player's proposal is complementary. For example, if the user proposes taking 2 scissors and the fork, and the other player takes 1 scissors and the thimble, that's complementary. The user's score would be 7, and the other's depends on their values} \\
            }
        }
    }
    & \\ \\

    \theutterance \stepcounter{utterance}  
    & & & \multicolumn{4}{p{0.6\linewidth}}{
        \cellcolor[rgb]{0.9,0.9,0.9}{
            \makecell[{{p{\linewidth}}}]{
                \texttt{\tiny{[GM$\langle$P2]}}
                \texttt{<think>} \\
\texttt{Okay, let's try to figure out the best way to approach this. The user has to divide 3 scissors, 1 thimble, and 1 fork between themselves and another player. Their own values are scissors: 3, thimble: 0, fork: 1. The other player's values are secret, so the user needs to negotiate to maximize the total score.} \\
\\ 
\texttt{First, the user's highest value is on scissors, so they probably want as many as possible. But there are three scissors, so maybe the other player can take some too. The fork is worth 1 to the user, and the thimble is worth 0. The other player might value the thimble or the fork more. Since there's only one of each, it's important to allocate them where they're most valued.} \\
\\ 
\texttt{The user can send up to 5 messages. They need to communicate their preferences and find out the other player's. But since the other player's values are secret, the user has to make proposals and see if they're complementary.} \\
\\ 
\texttt{One approach is to propose taking 2 scissors and the fork, leaving 1 scissors and the thimble. That way, the user gets 2*3 +1 =7 points. If the other player values the thimble, they can take it and the remaining scissors. But if the other player values the fork more, maybe they'd want it instead. Alternatively, the user could take 1 scissors and the fork, leaving 2 scissors and the thimble. That would give the user 3 +1 =4, but maybe the other player can get more from the remaining items.} \\
\\ 
\texttt{Wait, the goal is to maximize the total score. So the user should aim for a split where both players get the items they value most. Since the user's thimble is 0, it's better to let the other player have it if they value it. The fork is worth 1 to the user, so if the other player values it more, it's better to let them have it. But how to find out?} \\
\\ 
\texttt{Maybe the user should propose different splits and see if the other player's proposal is complementary. For example, if the user proposes taking 2 scissors and the fork, and the other player takes 1 scissors and the thimble, that's complementary. The user's score would be 7, and the other's depends on their values} \\
            }
        }
    }
     \\ \\

    \theutterance \stepcounter{utterance}  
    & & & \multicolumn{2}{p{0.3\linewidth}}{
        \cellcolor[rgb]{0.9,0.9,0.9}{
            \makecell[{{p{\linewidth}}}]{
                \texttt{\tiny{[GM$|$GM]}}
                \texttt{continue} \\
            }
        }
    }
    & & \\ \\

    \theutterance \stepcounter{utterance}  
    & & & \multicolumn{2}{p{0.3\linewidth}}{
        \cellcolor[rgb]{0.9,0.9,0.9}{
            \makecell[{{p{\linewidth}}}]{
                \texttt{\tiny{[GM$|$GM]}}
                \texttt{<think>} \\
\texttt{Okay, let's try to figure out the best way to approach this. The user has to divide 3 scissors, 1 thimble, and 1 fork between themselves and another player. Their own values are scissors: 3, thimble: 0, fork: 1. The other player's values are secret, so the user needs to negotiate to maximize the total score.} \\
\\ 
\texttt{First, the user's highest value is on scissors, so they probably want as many as possible. But there are three scissors, so maybe the other player can take some too. The fork is worth 1 to the user, and the thimble is worth 0. The other player might value the thimble or the fork more. Since there's only one of each, it's important to allocate them where they're most valued.} \\
\\ 
\texttt{The user can send up to 5 messages. They need to communicate their preferences and find out the other player's. But since the other player's values are secret, the user has to make proposals and see if they're complementary.} \\
\\ 
\texttt{One approach is to propose taking 2 scissors and the fork, leaving 1 scissors and the thimble. That way, the user gets 2*3 +1 =7 points. If the other player values the thimble, they can take it and the remaining scissors. But if the other player values the fork more, maybe they'd want it instead. Alternatively, the user could take 1 scissors and the fork, leaving 2 scissors and the thimble. That would give the user 3 +1 =4, but maybe the other player can get more from the remaining items.} \\
\\ 
\texttt{Wait, the goal is to maximize the total score. So the user should aim for a split where both players get the items they value most. Since the user's thimble is 0, it's better to let the other player have it if they value it. The fork is worth 1 to the user, so if the other player values it more, it's better to let them have it. But how to find out?} \\
\\ 
\texttt{Maybe the user should propose different splits and see if the other player's proposal is complementary. For example, if the user proposes taking 2 scissors and the fork, and the other player takes 1 scissors and the thimble, that's complementary. The user's score would be 7, and the other's depends on their values} \\
            }
        }
    }
    & & \\ \\

    \theutterance \stepcounter{utterance}  
    & & \multicolumn{4}{p{0.6\linewidth}}{
        \cellcolor[rgb]{0.9,0.9,0.9}{
            \makecell[{{p{\linewidth}}}]{
                \texttt{\tiny{[P1$\langle$GM]}}
                \texttt{<think>} \\
\texttt{Okay, let's try to figure out the best way to approach this. The user has to divide 3 scissors, 1 thimble, and 1 fork between themselves and another player. Their own values are scissors: 3, thimble: 0, fork: 1. The other player's values are secret, so the user needs to negotiate to maximize the total score.} \\
\\ 
\texttt{First, the user's highest value is on scissors, so they probably want as many as possible. But there are three scissors, so maybe the other player can take some too. The fork is worth 1 to the user, and the thimble is worth 0. The other player might value the thimble or the fork more. Since there's only one of each, it's important to allocate them where they're most valued.} \\
\\ 
\texttt{The user can send up to 5 messages. They need to communicate their preferences and find out the other player's. But since the other player's values are secret, the user has to make proposals and see if they're complementary.} \\
\\ 
\texttt{One approach is to propose taking 2 scissors and the fork, leaving 1 scissors and the thimble. That way, the user gets 2*3 +1 =7 points. If the other player values the thimble, they can take it and the remaining scissors. But if the other player values the fork more, maybe they'd want it instead. Alternatively, the user could take 1 scissors and the fork, leaving 2 scissors and the thimble. That would give the user 3 +1 =4, but maybe the other player can get more from the remaining items.} \\
\\ 
\texttt{Wait, the goal is to maximize the total score. So the user should aim for a split where both players get the items they value most. Since the user's thimble is 0, it's better to let the other player have it if they value it. The fork is worth 1 to the user, so if the other player values it more, it's better to let them have it. But how to find out?} \\
\\ 
\texttt{Maybe the user should propose different splits and see if the other player's proposal is complementary. For example, if the user proposes taking 2 scissors and the fork, and the other player takes 1 scissors and the thimble, that's complementary. The user's score would be 7, and the other's depends on their values} \\
            }
        }
    }
    & \\ \\

    \theutterance \stepcounter{utterance}  
    & \multicolumn{4}{p{0.6\linewidth}}{
        \cellcolor[rgb]{0.9,0.9,0.9}{
            \makecell[{{p{\linewidth}}}]{
                \texttt{\tiny{[P1$\rangle$GM]}}
                \texttt{<think>} \\
\texttt{Okay, let's try to figure out the best way to approach this. The user has to divide 3 scissors, 1 thimble, and 1 fork between themselves and another player. Their own values are scissors: 3, thimble: 0, fork: 1. The other player's values are secret, so the user needs to negotiate to maximize the total score.} \\
\\ 
\texttt{First, the user's highest value is on scissors, so they probably want as many as possible. But there are three scissors, so maybe the other player can take some too. The fork is worth 1 to the user, and the thimble is worth 0. The other player might value the thimble or the fork more. Since there's only one of each, it's important to allocate them where they're most valued.} \\
\\ 
\texttt{The user can send up to 5 messages. They need to communicate their preferences and find out the other player's. But since the other player's values are secret, the user has to make proposals and see if they're complementary.} \\
\\ 
\texttt{One approach is to propose taking 2 scissors and the fork, leaving 1 scissors and the thimble. That way, the user gets 2*3 +1 =7 points. If the other player values the thimble, they can take it and the remaining scissors. But if the other player values the fork more, maybe they'd want it instead. Alternatively, the user could take 1 scissors and the fork, leaving 2 scissors and the thimble. That would give the user 3 +1 =4, but maybe the other player can get more from the remaining items.} \\
\\ 
\texttt{Wait, the goal is to maximize the total score. So the user should aim for a split where both players get the items they value most. Since the user's thimble is 0, it's better to let the other player have it if they value it. The fork is worth 1 to the user, so if the other player values it more, it's better to let them have it. But how to find out?} \\
\\ 
\texttt{Maybe the user should propose different splits and see if the other player's proposal is complementary. For example, if the user proposes taking 2 scissors and the fork, and the other player takes 1 scissors and the thimble, that's complementary. The user's score would be 7, and the other's depends on their values} \\
            }
        }
    }
    & & \\ \\

    \theutterance \stepcounter{utterance}  
    & & & \multicolumn{2}{p{0.3\linewidth}}{
        \cellcolor[rgb]{0.9,0.9,0.9}{
            \makecell[{{p{\linewidth}}}]{
                \texttt{\tiny{[GM$|$GM]}}
                \texttt{continue} \\
            }
        }
    }
    & & \\ \\

    \theutterance \stepcounter{utterance}  
    & & & \multicolumn{2}{p{0.3\linewidth}}{
        \cellcolor[rgb]{0.9,0.9,0.9}{
            \makecell[{{p{\linewidth}}}]{
                \texttt{\tiny{[GM$|$GM]}}
                \texttt{<think>} \\
\texttt{Okay, let's try to figure out the best way to approach this. The user has to divide 3 scissors, 1 thimble, and 1 fork between themselves and another player. Their own values are scissors: 3, thimble: 0, fork: 1. The other player's values are secret, so the user needs to negotiate to maximize the total score.} \\
\\ 
\texttt{First, the user's highest value is on scissors, so they probably want as many as possible. But there are three scissors, so maybe the other player can take some too. The fork is worth 1 to the user, and the thimble is worth 0. The other player might value the thimble or the fork more. Since there's only one of each, it's important to allocate them where they're most valued.} \\
\\ 
\texttt{The user can send up to 5 messages. They need to communicate their preferences and find out the other player's. But since the other player's values are secret, the user has to make proposals and see if they're complementary.} \\
\\ 
\texttt{One approach is to propose taking 2 scissors and the fork, leaving 1 scissors and the thimble. That way, the user gets 2*3 +1 =7 points. If the other player values the thimble, they can take it and the remaining scissors. But if the other player values the fork more, maybe they'd want it instead. Alternatively, the user could take 1 scissors and the fork, leaving 2 scissors and the thimble. That would give the user 3 +1 =4, but maybe the other player can get more from the remaining items.} \\
\\ 
\texttt{Wait, the goal is to maximize the total score. So the user should aim for a split where both players get the items they value most. Since the user's thimble is 0, it's better to let the other player have it if they value it. The fork is worth 1 to the user, so if the other player values it more, it's better to let them have it. But how to find out?} \\
\\ 
\texttt{Maybe the user should propose different splits and see if the other player's proposal is complementary. For example, if the user proposes taking 2 scissors and the fork, and the other player takes 1 scissors and the thimble, that's complementary. The user's score would be 7, and the other's depends on their values} \\
            }
        }
    }
    & & \\ \\

    \theutterance \stepcounter{utterance}  
    & & \multicolumn{4}{p{0.6\linewidth}}{
        \cellcolor[rgb]{0.9,0.9,0.9}{
            \makecell[{{p{\linewidth}}}]{
                \texttt{\tiny{[GM$\rangle$P2]}}
                \texttt{<think>} \\
\texttt{Okay, let's try to figure out the best way to approach this. The user has to divide 3 scissors, 1 thimble, and 1 fork between themselves and another player. Their own values are scissors: 3, thimble: 0, fork: 1. The other player's values are secret, so the user needs to negotiate to maximize the total score.} \\
\\ 
\texttt{First, the user's highest value is on scissors, so they probably want as many as possible. But there are three scissors, so maybe the other player can take some too. The fork is worth 1 to the user, and the thimble is worth 0. The other player might value the thimble or the fork more. Since there's only one of each, it's important to allocate them where they're most valued.} \\
\\ 
\texttt{The user can send up to 5 messages. They need to communicate their preferences and find out the other player's. But since the other player's values are secret, the user has to make proposals and see if they're complementary.} \\
\\ 
\texttt{One approach is to propose taking 2 scissors and the fork, leaving 1 scissors and the thimble. That way, the user gets 2*3 +1 =7 points. If the other player values the thimble, they can take it and the remaining scissors. But if the other player values the fork more, maybe they'd want it instead. Alternatively, the user could take 1 scissors and the fork, leaving 2 scissors and the thimble. That would give the user 3 +1 =4, but maybe the other player can get more from the remaining items.} \\
\\ 
\texttt{Wait, the goal is to maximize the total score. So the user should aim for a split where both players get the items they value most. Since the user's thimble is 0, it's better to let the other player have it if they value it. The fork is worth 1 to the user, so if the other player values it more, it's better to let them have it. But how to find out?} \\
\\ 
\texttt{Maybe the user should propose different splits and see if the other player's proposal is complementary. For example, if the user proposes taking 2 scissors and the fork, and the other player takes 1 scissors and the thimble, that's complementary. The user's score would be 7, and the other's depends on their values} \\
            }
        }
    }
    & \\ \\

    \theutterance \stepcounter{utterance}  
    & & & \multicolumn{4}{p{0.6\linewidth}}{
        \cellcolor[rgb]{0.9,0.9,0.9}{
            \makecell[{{p{\linewidth}}}]{
                \texttt{\tiny{[GM$\langle$P2]}}
                \texttt{<think>} \\
\texttt{Okay, let's try to figure out the best way to approach this. The user has to divide 3 scissors, 1 thimble, and 1 fork between themselves and another player. Their own values are scissors: 3, thimble: 0, fork: 1. The other player's values are secret, so the user needs to negotiate to maximize the total score.} \\
\\ 
\texttt{First, the user's highest value is on scissors, so they probably want as many as possible. But there are three scissors, so maybe the other player can take some too. The fork is worth 1 to the user, and the thimble is worth 0. The other player might value the thimble or the fork more. Since there's only one of each, it's important to allocate them where they're most valued.} \\
\\ 
\texttt{The user can send up to 5 messages. They need to communicate their preferences and find out the other player's. But since the other player's values are secret, the user has to make proposals and see if they're complementary.} \\
\\ 
\texttt{One approach is to propose taking 2 scissors and the fork, leaving 1 scissors and the thimble. That way, the user gets 2*3 +1 =7 points. If the other player values the thimble, they can take it and the remaining scissors. But if the other player values the fork more, maybe they'd want it instead. Alternatively, the user could take 1 scissors and the fork, leaving 2 scissors and the thimble. That would give the user 3 +1 =4, but maybe the other player can get more from the remaining items.} \\
\\ 
\texttt{Wait, the goal is to maximize the total score. So the user should aim for a split where both players get the items they value most. Since the user's thimble is 0, it's better to let the other player have it if they value it. The fork is worth 1 to the user, so if the other player values it more, it's better to let them have it. But how to find out?} \\
\\ 
\texttt{Maybe the user should propose different splits and see if the other player's proposal is complementary. For example, if the user proposes taking 2 scissors and the fork, and the other player takes 1 scissors and the thimble, that's complementary. The user's score would be 7, and the other's depends on their values} \\
            }
        }
    }
     \\ \\

    \theutterance \stepcounter{utterance}  
    & & & \multicolumn{2}{p{0.3\linewidth}}{
        \cellcolor[rgb]{0.9,0.9,0.9}{
            \makecell[{{p{\linewidth}}}]{
                \texttt{\tiny{[GM$|$GM]}}
                \texttt{continue} \\
            }
        }
    }
    & & \\ \\

    \theutterance \stepcounter{utterance}  
    & & & \multicolumn{2}{p{0.3\linewidth}}{
        \cellcolor[rgb]{0.9,0.9,0.9}{
            \makecell[{{p{\linewidth}}}]{
                \texttt{\tiny{[GM$|$GM]}}
                \texttt{<think>} \\
\texttt{Okay, let's try to figure out the best way to approach this. The user has to divide 3 scissors, 1 thimble, and 1 fork between themselves and another player. Their own values are scissors: 3, thimble: 0, fork: 1. The other player's values are secret, so the user needs to negotiate to maximize the total score.} \\
\\ 
\texttt{First, the user's highest value is on scissors, so they probably want as many as possible. But there are three scissors, so maybe the other player can take some too. The fork is worth 1 to the user, and the thimble is worth 0. The other player might value the thimble or the fork more. Since there's only one of each, it's important to allocate them where they're most valued.} \\
\\ 
\texttt{The user can send up to 5 messages. They need to communicate their preferences and find out the other player's. But since the other player's values are secret, the user has to make proposals and see if they're complementary.} \\
\\ 
\texttt{One approach is to propose taking 2 scissors and the fork, leaving 1 scissors and the thimble. That way, the user gets 2*3 +1 =7 points. If the other player values the thimble, they can take it and the remaining scissors. But if the other player values the fork more, maybe they'd want it instead. Alternatively, the user could take 1 scissors and the fork, leaving 2 scissors and the thimble. That would give the user 3 +1 =4, but maybe the other player can get more from the remaining items.} \\
\\ 
\texttt{Wait, the goal is to maximize the total score. So the user should aim for a split where both players get the items they value most. Since the user's thimble is 0, it's better to let the other player have it if they value it. The fork is worth 1 to the user, so if the other player values it more, it's better to let them have it. But how to find out?} \\
\\ 
\texttt{Maybe the user should propose different splits and see if the other player's proposal is complementary. For example, if the user proposes taking 2 scissors and the fork, and the other player takes 1 scissors and the thimble, that's complementary. The user's score would be 7, and the other's depends on their values} \\
            }
        }
    }
    & & \\ \\

    \theutterance \stepcounter{utterance}  
    & & \multicolumn{4}{p{0.6\linewidth}}{
        \cellcolor[rgb]{0.9,0.9,0.9}{
            \makecell[{{p{\linewidth}}}]{
                \texttt{\tiny{[P1$\langle$GM]}}
                \texttt{<think>} \\
\texttt{Okay, let's try to figure out the best way to approach this. The user has to divide 3 scissors, 1 thimble, and 1 fork between themselves and another player. Their own values are scissors: 3, thimble: 0, fork: 1. The other player's values are secret, so the user needs to negotiate to maximize the total score.} \\
\\ 
\texttt{First, the user's highest value is on scissors, so they probably want as many as possible. But there are three scissors, so maybe the other player can take some too. The fork is worth 1 to the user, and the thimble is worth 0. The other player might value the thimble or the fork more. Since there's only one of each, it's important to allocate them where they're most valued.} \\
\\ 
\texttt{The user can send up to 5 messages. They need to communicate their preferences and find out the other player's. But since the other player's values are secret, the user has to make proposals and see if they're complementary.} \\
\\ 
\texttt{One approach is to propose taking 2 scissors and the fork, leaving 1 scissors and the thimble. That way, the user gets 2*3 +1 =7 points. If the other player values the thimble, they can take it and the remaining scissors. But if the other player values the fork more, maybe they'd want it instead. Alternatively, the user could take 1 scissors and the fork, leaving 2 scissors and the thimble. That would give the user 3 +1 =4, but maybe the other player can get more from the remaining items.} \\
\\ 
\texttt{Wait, the goal is to maximize the total score. So the user should aim for a split where both players get the items they value most. Since the user's thimble is 0, it's better to let the other player have it if they value it. The fork is worth 1 to the user, so if the other player values it more, it's better to let them have it. But how to find out?} \\
\\ 
\texttt{Maybe the user should propose different splits and see if the other player's proposal is complementary. For example, if the user proposes taking 2 scissors and the fork, and the other player takes 1 scissors and the thimble, that's complementary. The user's score would be 7, and the other's depends on their values} \\
            }
        }
    }
    & \\ \\

    \theutterance \stepcounter{utterance}  
    & \multicolumn{4}{p{0.6\linewidth}}{
        \cellcolor[rgb]{0.9,0.9,0.9}{
            \makecell[{{p{\linewidth}}}]{
                \texttt{\tiny{[P1$\rangle$GM]}}
                \texttt{<think>} \\
\texttt{Okay, let's try to figure out the best way to approach this. The user has to divide 3 scissors, 1 thimble, and 1 fork between themselves and another player. Their own values are scissors: 3, thimble: 0, fork: 1. The other player's values are secret, so the user needs to negotiate to maximize the total score.} \\
\\ 
\texttt{First, the user's highest value is on scissors, so they probably want as many as possible. But there are three scissors, so maybe the other player can take some too. The fork is worth 1 to the user, and the thimble is worth 0. The other player might value the thimble or the fork more. Since there's only one of each, it's important to allocate them where they're most valued.} \\
\\ 
\texttt{The user can send up to 5 messages. They need to communicate their preferences and find out the other player's. But since the other player's values are secret, the user has to make proposals and see if they're complementary.} \\
\\ 
\texttt{One approach is to propose taking 2 scissors and the fork, leaving 1 scissors and the thimble. That way, the user gets 2*3 +1 =7 points. If the other player values the thimble, they can take it and the remaining scissors. But if the other player values the fork more, maybe they'd want it instead. Alternatively, the user could take 1 scissors and the fork, leaving 2 scissors and the thimble. That would give the user 3 +1 =4, but maybe the other player can get more from the remaining items.} \\
\\ 
\texttt{Wait, the goal is to maximize the total score. So the user should aim for a split where both players get the items they value most. Since the user's thimble is 0, it's better to let the other player have it if they value it. The fork is worth 1 to the user, so if the other player values it more, it's better to let them have it. But how to find out?} \\
\\ 
\texttt{Maybe the user should propose different splits and see if the other player's proposal is complementary. For example, if the user proposes taking 2 scissors and the fork, and the other player takes 1 scissors and the thimble, that's complementary. The user's score would be 7, and the other's depends on their values} \\
            }
        }
    }
    & & \\ \\

    \theutterance \stepcounter{utterance}  
    & & & \multicolumn{2}{p{0.3\linewidth}}{
        \cellcolor[rgb]{0.9,0.9,0.9}{
            \makecell[{{p{\linewidth}}}]{
                \texttt{\tiny{[GM$|$GM]}}
                \texttt{continue} \\
            }
        }
    }
    & & \\ \\

    \theutterance \stepcounter{utterance}  
    & & & \multicolumn{2}{p{0.3\linewidth}}{
        \cellcolor[rgb]{0.9,0.9,0.9}{
            \makecell[{{p{\linewidth}}}]{
                \texttt{\tiny{[GM$|$GM]}}
                \texttt{<think>} \\
\texttt{Okay, let's try to figure out the best way to approach this. The user has to divide 3 scissors, 1 thimble, and 1 fork between themselves and another player. Their own values are scissors: 3, thimble: 0, fork: 1. The other player's values are secret, so the user needs to negotiate to maximize the total score.} \\
\\ 
\texttt{First, the user's highest value is on scissors, so they probably want as many as possible. But there are three scissors, so maybe the other player can take some too. The fork is worth 1 to the user, and the thimble is worth 0. The other player might value the thimble or the fork more. Since there's only one of each, it's important to allocate them where they're most valued.} \\
\\ 
\texttt{The user can send up to 5 messages. They need to communicate their preferences and find out the other player's. But since the other player's values are secret, the user has to make proposals and see if they're complementary.} \\
\\ 
\texttt{One approach is to propose taking 2 scissors and the fork, leaving 1 scissors and the thimble. That way, the user gets 2*3 +1 =7 points. If the other player values the thimble, they can take it and the remaining scissors. But if the other player values the fork more, maybe they'd want it instead. Alternatively, the user could take 1 scissors and the fork, leaving 2 scissors and the thimble. That would give the user 3 +1 =4, but maybe the other player can get more from the remaining items.} \\
\\ 
\texttt{Wait, the goal is to maximize the total score. So the user should aim for a split where both players get the items they value most. Since the user's thimble is 0, it's better to let the other player have it if they value it. The fork is worth 1 to the user, so if the other player values it more, it's better to let them have it. But how to find out?} \\
\\ 
\texttt{Maybe the user should propose different splits and see if the other player's proposal is complementary. For example, if the user proposes taking 2 scissors and the fork, and the other player takes 1 scissors and the thimble, that's complementary. The user's score would be 7, and the other's depends on their values} \\
            }
        }
    }
    & & \\ \\

    \theutterance \stepcounter{utterance}  
    & & \multicolumn{4}{p{0.6\linewidth}}{
        \cellcolor[rgb]{0.9,0.9,0.9}{
            \makecell[{{p{\linewidth}}}]{
                \texttt{\tiny{[GM$\rangle$P2]}}
                \texttt{<think>} \\
\texttt{Okay, let's try to figure out the best way to approach this. The user has to divide 3 scissors, 1 thimble, and 1 fork between themselves and another player. Their own values are scissors: 3, thimble: 0, fork: 1. The other player's values are secret, so the user needs to negotiate to maximize the total score.} \\
\\ 
\texttt{First, the user's highest value is on scissors, so they probably want as many as possible. But there are three scissors, so maybe the other player can take some too. The fork is worth 1 to the user, and the thimble is worth 0. The other player might value the thimble or the fork more. Since there's only one of each, it's important to allocate them where they're most valued.} \\
\\ 
\texttt{The user can send up to 5 messages. They need to communicate their preferences and find out the other player's. But since the other player's values are secret, the user has to make proposals and see if they're complementary.} \\
\\ 
\texttt{One approach is to propose taking 2 scissors and the fork, leaving 1 scissors and the thimble. That way, the user gets 2*3 +1 =7 points. If the other player values the thimble, they can take it and the remaining scissors. But if the other player values the fork more, maybe they'd want it instead. Alternatively, the user could take 1 scissors and the fork, leaving 2 scissors and the thimble. That would give the user 3 +1 =4, but maybe the other player can get more from the remaining items.} \\
\\ 
\texttt{Wait, the goal is to maximize the total score. So the user should aim for a split where both players get the items they value most. Since the user's thimble is 0, it's better to let the other player have it if they value it. The fork is worth 1 to the user, so if the other player values it more, it's better to let them have it. But how to find out?} \\
\\ 
\texttt{Maybe the user should propose different splits and see if the other player's proposal is complementary. For example, if the user proposes taking 2 scissors and the fork, and the other player takes 1 scissors and the thimble, that's complementary. The user's score would be 7, and the other's depends on their values} \\
            }
        }
    }
    & \\ \\

    \theutterance \stepcounter{utterance}  
    & & & \multicolumn{4}{p{0.6\linewidth}}{
        \cellcolor[rgb]{0.9,0.9,0.9}{
            \makecell[{{p{\linewidth}}}]{
                \texttt{\tiny{[GM$\langle$P2]}}
                \texttt{<think>} \\
\texttt{Okay, let's try to figure out the best way to approach this. The user has to divide 3 scissors, 1 thimble, and 1 fork between themselves and another player. Their own values are scissors: 3, thimble: 0, fork: 1. The other player's values are secret, so the user needs to negotiate to maximize the total score.} \\
\\ 
\texttt{First, the user's highest value is on scissors, so they probably want as many as possible. But there are three scissors, so maybe the other player can take some too. The fork is worth 1 to the user, and the thimble is worth 0. The other player might value the thimble or the fork more. Since there's only one of each, it's important to allocate them where they're most valued.} \\
\\ 
\texttt{The user can send up to 5 messages. They need to communicate their preferences and find out the other player's. But since the other player's values are secret, the user has to make proposals and see if they're complementary.} \\
\\ 
\texttt{One approach is to propose taking 2 scissors and the fork, leaving 1 scissors and the thimble. That way, the user gets 2*3 +1 =7 points. If the other player values the thimble, they can take it and the remaining scissors. But if the other player values the fork more, maybe they'd want it instead. Alternatively, the user could take 1 scissors and the fork, leaving 2 scissors and the thimble. That would give the user 3 +1 =4, but maybe the other player can get more from the remaining items.} \\
\\ 
\texttt{Wait, the goal is to maximize the total score. So the user should aim for a split where both players get the items they value most. Since the user's thimble is 0, it's better to let the other player have it if they value it. The fork is worth 1 to the user, so if the other player values it more, it's better to let them have it. But how to find out?} \\
\\ 
\texttt{Maybe the user should propose different splits and see if the other player's proposal is complementary. For example, if the user proposes taking 2 scissors and the fork, and the other player takes 1 scissors and the thimble, that's complementary. The user's score would be 7, and the other's depends on their values} \\
            }
        }
    }
     \\ \\

    \theutterance \stepcounter{utterance}  
    & & & \multicolumn{2}{p{0.3\linewidth}}{
        \cellcolor[rgb]{0.9,0.9,0.9}{
            \makecell[{{p{\linewidth}}}]{
                \texttt{\tiny{[GM$|$GM]}}
                \texttt{continue} \\
            }
        }
    }
    & & \\ \\

    \theutterance \stepcounter{utterance}  
    & & & \multicolumn{2}{p{0.3\linewidth}}{
        \cellcolor[rgb]{0.9,0.9,0.9}{
            \makecell[{{p{\linewidth}}}]{
                \texttt{\tiny{[GM$|$GM]}}
                \texttt{<think>} \\
\texttt{Okay, let's try to figure out the best way to approach this. The user has to divide 3 scissors, 1 thimble, and 1 fork between themselves and another player. Their own values are scissors: 3, thimble: 0, fork: 1. The other player's values are secret, so the user needs to negotiate to maximize the total score.} \\
\\ 
\texttt{First, the user's highest value is on scissors, so they probably want as many as possible. But there are three scissors, so maybe the other player can take some too. The fork is worth 1 to the user, and the thimble is worth 0. The other player might value the thimble or the fork more. Since there's only one of each, it's important to allocate them where they're most valued.} \\
\\ 
\texttt{The user can send up to 5 messages. They need to communicate their preferences and find out the other player's. But since the other player's values are secret, the user has to make proposals and see if they're complementary.} \\
\\ 
\texttt{One approach is to propose taking 2 scissors and the fork, leaving 1 scissors and the thimble. That way, the user gets 2*3 +1 =7 points. If the other player values the thimble, they can take it and the remaining scissors. But if the other player values the fork more, maybe they'd want it instead. Alternatively, the user could take 1 scissors and the fork, leaving 2 scissors and the thimble. That would give the user 3 +1 =4, but maybe the other player can get more from the remaining items.} \\
\\ 
\texttt{Wait, the goal is to maximize the total score. So the user should aim for a split where both players get the items they value most. Since the user's thimble is 0, it's better to let the other player have it if they value it. The fork is worth 1 to the user, so if the other player values it more, it's better to let them have it. But how to find out?} \\
\\ 
\texttt{Maybe the user should propose different splits and see if the other player's proposal is complementary. For example, if the user proposes taking 2 scissors and the fork, and the other player takes 1 scissors and the thimble, that's complementary. The user's score would be 7, and the other's depends on their values} \\
            }
        }
    }
    & & \\ \\

    \theutterance \stepcounter{utterance}  
    & & \multicolumn{4}{p{0.6\linewidth}}{
        \cellcolor[rgb]{0.9,0.9,0.9}{
            \makecell[{{p{\linewidth}}}]{
                \texttt{\tiny{[P1$\langle$GM]}}
                \texttt{<think>} \\
\texttt{Okay, let's try to figure out the best way to approach this. The user has to divide 3 scissors, 1 thimble, and 1 fork between themselves and another player. Their own values are scissors: 3, thimble: 0, fork: 1. The other player's values are secret, so the user needs to negotiate to maximize the total score.} \\
\\ 
\texttt{First, the user's highest value is on scissors, so they probably want as many as possible. But there are three scissors, so maybe the other player can take some too. The fork is worth 1 to the user, and the thimble is worth 0. The other player might value the thimble or the fork more. Since there's only one of each, it's important to allocate them where they're most valued.} \\
\\ 
\texttt{The user can send up to 5 messages. They need to communicate their preferences and find out the other player's. But since the other player's values are secret, the user has to make proposals and see if they're complementary.} \\
\\ 
\texttt{One approach is to propose taking 2 scissors and the fork, leaving 1 scissors and the thimble. That way, the user gets 2*3 +1 =7 points. If the other player values the thimble, they can take it and the remaining scissors. But if the other player values the fork more, maybe they'd want it instead. Alternatively, the user could take 1 scissors and the fork, leaving 2 scissors and the thimble. That would give the user 3 +1 =4, but maybe the other player can get more from the remaining items.} \\
\\ 
\texttt{Wait, the goal is to maximize the total score. So the user should aim for a split where both players get the items they value most. Since the user's thimble is 0, it's better to let the other player have it if they value it. The fork is worth 1 to the user, so if the other player values it more, it's better to let them have it. But how to find out?} \\
\\ 
\texttt{Maybe the user should propose different splits and see if the other player's proposal is complementary. For example, if the user proposes taking 2 scissors and the fork, and the other player takes 1 scissors and the thimble, that's complementary. The user's score would be 7, and the other's depends on their values} \\
\\ 
\\ 
\texttt{The time is up. It is now your turn to submit a secret proposal.} \\
            }
        }
    }
    & \\ \\

    \theutterance \stepcounter{utterance}  
    & \multicolumn{4}{p{0.6\linewidth}}{
        \cellcolor[rgb]{0.9,0.9,0.9}{
            \makecell[{{p{\linewidth}}}]{
                \texttt{\tiny{[P1$\rangle$GM]}}
                \texttt{<think>} \\
\texttt{Okay, let's try to figure out the best way to approach this. The user has to divide 3 scissors, 1 thimble, and 1 fork between themselves and another player. Their own values are scissors: 3, thimble: 0, fork: 1. The other player's values are secret, so the user needs to negotiate to maximize the total score.} \\
\\ 
\texttt{First, the user's highest value is on scissors, so they probably want as many as possible. But there are three scissors, so maybe the other player can take some too. The fork is worth 1 to the user, and the thimble is worth 0. The other player might value the thimble or the fork more. Since there's only one of each, it's important to allocate them where they're most valued.} \\
\\ 
\texttt{The user can send up to 5 messages. They need to communicate their preferences and find out the other player's. But since the other player's values are secret, the user has to make proposals and see if they're complementary.} \\
\\ 
\texttt{One approach is to propose taking 2 scissors and the fork, leaving 1 scissors and the thimble. That way, the user gets 2*3 +1 =7 points. If the other player values the thimble, they can take it and the remaining scissors. But if the other player values the fork more, maybe they'd want it instead. Alternatively, the user could take 1 scissors and the fork, leaving 2 scissors and the thimble. That would give the user 3 +1 =4, but maybe the other player can get more from the remaining items.} \\
\\ 
\texttt{Wait, the goal is to maximize the total score. So the user should aim for a split where both players get the items they value most. Since the user's thimble is 0, it's better to let the other player have it if they value it. The fork is worth 1 to the user, so if the other player values it more, it's better to let them have it. But how to find out?} \\
\\ 
\texttt{Maybe the user should propose different splits and see if the other player's proposal is complementary. For example, if the user proposes taking 2 scissors and the fork, and the other player takes 1 scissors and the thimble, that's complementary. The user's score would be 7, and the other's depends on their values} \\
            }
        }
    }
    & & \\ \\

    \theutterance \stepcounter{utterance}  
    & & & \multicolumn{2}{p{0.3\linewidth}}{
        \cellcolor[rgb]{0.9,0.9,0.9}{
            \makecell[{{p{\linewidth}}}]{
                \texttt{\tiny{[GM$|$GM]}}
                \texttt{continue} \\
            }
        }
    }
    & & \\ \\

    \theutterance \stepcounter{utterance}  
    & & & \multicolumn{2}{p{0.3\linewidth}}{
        \cellcolor[rgb]{0.9,0.9,0.9}{
            \makecell[{{p{\linewidth}}}]{
                \texttt{\tiny{[GM$|$GM]}}
                \texttt{player was instructed to make a proposal} \\
            }
        }
    }
    & & \\ \\

\end{supertabular}
}

\end{document}
